\documentclass[a4paper,11pt]{article}
\usepackage[utf8]{inputenc}
\usepackage[T1]{fontenc}
\usepackage[french]{babel}
\usepackage{amsmath,amssymb}
\usepackage{geometry}
\usepackage{tikz}
\usepackage{enumitem}
\usepackage{fancyhdr}
\usepackage{tcolorbox}
\usepackage{pgfplots}
\pgfplotsset{compat=1.18}
\usepackage{xcolor}

\geometry{margin=2cm}

\pagestyle{fancy}
\fancyhf{}
\lhead{BTS -- Mathématiques}
\rhead{Séries de Fourier -- \textcolor{red}{\textbf{CORRIGÉ}}}
\cfoot{\thepage}
\setlength{\headheight}{14pt}

\tcbset{
  formulebox/.style={colback=blue!5!white, colframe=blue!50!black, fonttitle=\bfseries, title=#1},
  exobox/.style={colback=orange!5!white, colframe=orange!60!black, fonttitle=\bfseries, title=#1},
  corrigebox/.style={colback=green!5!white, colframe=green!50!black, sharp corners, boxrule=0.5pt, left=4pt, right=4pt, top=4pt, bottom=4pt}
}

\begin{document}

\begin{center}
  {\LARGE \textbf{Activité -- Séries de Fourier}} \\[0.3cm]
  {\Large Calcul de $a_0$ : la valeur moyenne d'un signal périodique} \\[0.2cm]
  \textcolor{red}{\Large \textbf{--- CORRIGÉ ---}} \\[0.2cm]
  \rule{0.6\textwidth}{0.4pt}
\end{center}

\vspace{0.5cm}

%% --- RAPPEL DE COURS ---
\begin{tcolorbox}[formulebox={Rappel : valeur moyenne d'un signal périodique}]
  Soit $f$ une fonction périodique de période $T$. Le coefficient $a_0$ représente la \textbf{valeur moyenne} de $f$ sur une période :
  \[
    \boxed{a_0 = \frac{1}{T} \int_0^{T} f(t)\, dt}
  \]
  \textbf{Interprétation physique :} $a_0$ correspond à la \textbf{composante continue} du signal, c'est-à-dire la valeur autour de laquelle le signal oscille.
\end{tcolorbox}

\vspace{0.5cm}

%% ================================================================
%% EXERCICE 1
%% ================================================================
\begin{tcolorbox}[exobox={Exercice 1 -- Signal créneau}]
  On considère le signal périodique $f$ de période $T = 2$ défini sur $[0\,;\,2[$ par :
  \[
    f(t) = \begin{cases} 3 & \text{si } 0 \leq t < 1 \\ -1 & \text{si } 1 \leq t < 2 \end{cases}
  \]
\end{tcolorbox}

\begin{center}
  \begin{tikzpicture}
    \begin{axis}[
      axis lines=middle,
      xlabel={$t$},
      ylabel={$f(t)$},
      xmin=-2.5, xmax=4.5,
      ymin=-2, ymax=4,
      xtick={-2,-1,0,1,2,3,4},
      ytick={-1,0,1,2,3},
      width=12cm, height=5cm,
      grid=major,
      grid style={dashed, gray!30},
    ]
    \addplot[blue, ultra thick, domain=-2:-1] {3};
    \addplot[blue, ultra thick, domain=-1:0] {-1};
    \addplot[blue, ultra thick, domain=0:1] {3};
    \addplot[blue, ultra thick, domain=1:2] {-1};
    \addplot[blue, ultra thick, domain=2:3] {3};
    \addplot[blue, ultra thick, domain=3:4] {-1};
    \addplot[blue, dashed, thin] coordinates {(1,3)(1,-1)};
    \addplot[blue, dashed, thin] coordinates {(3,3)(3,-1)};
    \addplot[blue, dashed, thin] coordinates {(-1,3)(-1,-1)};
    % Droite a0
    \addplot[red, dashed, thick, domain=-2.5:4.5] {1};
    \node[red] at (axis cs:4.2,1.4) {$a_0=1$};
    \end{axis}
  \end{tikzpicture}
\end{center}

\begin{enumerate}[label=\textbf{\alph*)}]
  \item Rappeler la valeur de la période $T$ et identifier les bornes d'intégration.
  \begin{tcolorbox}[corrigebox]
    La période est $T = 2$. On intègre de $0$ à $T = 2$.
  \end{tcolorbox}

  \item Écrire l'intégrale permettant de calculer $a_0$ en la décomposant selon les deux intervalles.
  \begin{tcolorbox}[corrigebox]
    La fonction $f$ est définie par morceaux, on décompose l'intégrale :
    \[
      a_0 = \frac{1}{T}\int_0^T f(t)\,dt = \frac{1}{2}\left[\int_0^1 3\,dt + \int_1^2 (-1)\,dt\right]
    \]
  \end{tcolorbox}

  \item Calculer $a_0$.
  \begin{tcolorbox}[corrigebox]
    \begin{align*}
      a_0 &= \frac{1}{2}\left[\Big[3t\Big]_0^1 + \Big[-t\Big]_1^2\right] \\
          &= \frac{1}{2}\left[(3\times 1 - 3\times 0) + (-2 - (-1))\right] \\
          &= \frac{1}{2}\left[3 + (-1)\right] \\
          &= \frac{1}{2}\times 2
    \end{align*}
    \[\boxed{a_0 = 1}\]
  \end{tcolorbox}

  \item Tracer en pointillé rouge la droite $y = a_0$ sur le graphique ci-dessus. Le résultat est-il cohérent graphiquement ?
  \begin{tcolorbox}[corrigebox]
    La droite $y = 1$ est tracée en rouge sur le graphique. Le résultat est cohérent : le signal passe la moitié du temps à $3$ et l'autre moitié à $-1$. La valeur $1$ se situe bien entre ces deux niveaux, et l'aire au-dessus de la droite (entre $y=1$ et $y=3$) compense l'aire en dessous (entre $y=-1$ et $y=1$).
  \end{tcolorbox}
\end{enumerate}

\vspace{0.3cm}

%% ================================================================
%% EXERCICE 2
%% ================================================================
\begin{tcolorbox}[exobox={Exercice 2 -- Signal en dents de scie}]
  On considère le signal périodique $g$ de période $T = 4$ défini sur $[0\,;\,4[$ par :
  \[
    g(t) = \frac{3}{4}\,t
  \]
\end{tcolorbox}

\begin{center}
  \begin{tikzpicture}
    \begin{axis}[
      axis lines=middle,
      xlabel={$t$},
      ylabel={$g(t)$},
      xmin=-4.5, xmax=8.5,
      ymin=-0.5, ymax=4,
      xtick={-4,-2,0,2,4,6,8},
      ytick={0,1,1.5,2,3},
      width=12cm, height=5cm,
      grid=major,
      grid style={dashed, gray!30},
    ]
    \addplot[red, ultra thick, domain=-4:0] {0.75*(x+4)};
    \addplot[red, ultra thick, domain=0:4] {0.75*x};
    \addplot[red, ultra thick, domain=4:8] {0.75*(x-4)};
    \addplot[red, dashed, thin] coordinates {(0,3)(0,0)};
    \addplot[red, dashed, thin] coordinates {(4,3)(4,0)};
    % Droite a0
    \addplot[magenta, dashed, thick, domain=-4.5:8.5] {1.5};
    \node[magenta] at (axis cs:7.5,1.85) {$a_0=\frac{3}{2}$};
    \end{axis}
  \end{tikzpicture}
\end{center}

\begin{enumerate}[label=\textbf{\alph*)}]
  \item Deviner graphiquement la valeur de $a_0$ avant tout calcul. Justifier.
  \begin{tcolorbox}[corrigebox]
    Sur une période, le signal est une droite qui va de $0$ à $3$. Par symétrie, la valeur moyenne se situe à mi-hauteur :
    \[
      a_0 = \frac{0+3}{2} = \frac{3}{2} = 1{,}5
    \]
  \end{tcolorbox}

  \item Écrire l'intégrale permettant de calculer $a_0$.
  \begin{tcolorbox}[corrigebox]
    \[
      a_0 = \frac{1}{T}\int_0^T g(t)\,dt = \frac{1}{4}\int_0^4 \frac{3}{4}\,t\,dt
    \]
  \end{tcolorbox}

  \item Calculer $a_0$. Le résultat confirme-t-il votre intuition graphique ?
  \begin{tcolorbox}[corrigebox]
    \begin{align*}
      a_0 &= \frac{1}{4}\int_0^4 \frac{3}{4}\,t\,dt = \frac{1}{4}\times\frac{3}{4}\int_0^4 t\,dt \\
          &= \frac{3}{16}\left[\frac{t^2}{2}\right]_0^4 = \frac{3}{16}\times\frac{16}{2} = \frac{3}{16}\times 8
    \end{align*}
    \[\boxed{a_0 = \frac{3}{2} = 1{,}5}\]
    Le résultat confirme bien l'intuition graphique.
  \end{tcolorbox}
\end{enumerate}

\vspace{0.3cm}

%% ================================================================
%% EXERCICE 3
%% ================================================================
\begin{tcolorbox}[exobox={Exercice 3 -- Signal sinusoïdal redressé}]
  On considère le signal périodique $h$ de période $T = \pi$ défini par :
  \[
    h(t) = \left| \sin(t) \right|
  \]
\end{tcolorbox}

\begin{center}
  \begin{tikzpicture}
    \begin{axis}[
      axis lines=middle,
      xlabel={$t$},
      ylabel={$h(t)$},
      xmin=-3.5, xmax=10,
      ymin=-0.3, ymax=1.5,
      xtick={-3.14159,0,3.14159,6.28318,9.42478},
      xticklabels={$-\pi$,$0$,$\pi$,$2\pi$,$3\pi$},
      ytick={0,0.5,0.6366,1},
      yticklabels={$0$,$0{,}5$,$a_0$,$1$},
      width=12cm, height=5cm,
      grid=major,
      grid style={dashed, gray!30},
      samples=200,
    ]
    \addplot[green!60!black, ultra thick, domain=-3.14159:9.5] {abs(sin(deg(x)))};
    \addplot[magenta, dashed, thick, domain=-3.5:10] {0.6366};
    \end{axis}
  \end{tikzpicture}
\end{center}

\begin{enumerate}[label=\textbf{\alph*)}]
  \item Justifier que la période de $h$ est bien $T = \pi$.
  \begin{tcolorbox}[corrigebox]
    La fonction $\sin(t)$ a pour période $2\pi$. Cependant, en prenant la valeur absolue, les arches négatives sont « repliées » vers le haut et deviennent identiques aux arches positives. La fonction $|\sin(t)|$ se répète donc tous les $\pi$ :
    \[
      |\sin(t+\pi)| = |-\sin(t)| = |\sin(t)|
    \]
    Donc $T = \pi$.
  \end{tcolorbox}

  \item Simplifier $|{\sin(t)}|$ sur l'intervalle $[0\,;\,\pi]$.
  \begin{tcolorbox}[corrigebox]
    Sur $[0\,;\,\pi]$, $\sin(t) \geq 0$, donc $|\sin(t)| = \sin(t)$.
  \end{tcolorbox}

  \item Écrire puis calculer l'intégrale donnant $a_0$.
  \begin{tcolorbox}[corrigebox]
    \begin{align*}
      a_0 &= \frac{1}{T}\int_0^T h(t)\,dt = \frac{1}{\pi}\int_0^{\pi} \sin(t)\,dt \\
          &= \frac{1}{\pi}\Big[-\cos(t)\Big]_0^{\pi} \\
          &= \frac{1}{\pi}\left[-\cos(\pi) - \bigl(-\cos(0)\bigr)\right] \\
          &= \frac{1}{\pi}\left[-(-1) + 1\right] = \frac{1}{\pi}\times 2
    \end{align*}
    \[\boxed{a_0 = \frac{2}{\pi}}\]
  \end{tcolorbox}

  \item Donner la valeur exacte de $a_0$ puis une valeur approchée à $10^{-2}$ près.
  \begin{tcolorbox}[corrigebox]
    Valeur exacte : $a_0 = \dfrac{2}{\pi}$

    \medskip
    Valeur approchée : $a_0 \approx \dfrac{2}{3{,}1416} \approx 0{,}64$
  \end{tcolorbox}
\end{enumerate}

\vspace{0.3cm}

%% ================================================================
%% EXERCICE 4
%% ================================================================
\begin{tcolorbox}[exobox={Exercice 4 -- Synthèse}]
  On considère le signal périodique $s$ de période $T = 6$ défini sur $[0\,;\,6[$ par :
  \[
    s(t) = \begin{cases} 2t & \text{si } 0 \leq t < 3 \\ 0 & \text{si } 3 \leq t < 6 \end{cases}
  \]
\end{tcolorbox}

\begin{enumerate}[label=\textbf{\alph*)}]
  \item Tracer le signal $s$ sur l'intervalle $[-6\,;\,12]$ dans le repère ci-dessous.

  \begin{center}
    \begin{tikzpicture}
      \begin{axis}[
        axis lines=middle,
        xlabel={$t$},
        ylabel={$s(t)$},
        xmin=-6.5, xmax=12.5,
        ymin=-1, ymax=7,
        xtick={-6,-3,0,3,6,9,12},
        ytick={0,1,2,3,4,5,6},
        width=14cm, height=6cm,
        grid=major,
        grid style={dashed, gray!30},
      ]
      % Période [-6, 0[
      \addplot[blue, ultra thick, domain=-6:-3] {2*(x+6)};
      \addplot[blue, ultra thick, domain=-3:0] {0};
      % Période [0, 6[
      \addplot[blue, ultra thick, domain=0:3] {2*x};
      \addplot[blue, ultra thick, domain=3:6] {0};
      % Période [6, 12[
      \addplot[blue, ultra thick, domain=6:9] {2*(x-6)};
      \addplot[blue, ultra thick, domain=9:12] {0};
      % Discontinuités
      \addplot[blue, dashed, thin] coordinates {(3,6)(3,0)};
      \addplot[blue, dashed, thin] coordinates {(9,6)(9,0)};
      \addplot[blue, dashed, thin] coordinates {(-3,6)(-3,0)};
      % Droite a0
      \addplot[red, dashed, thick, domain=-6.5:12.5] {3};
      \node[red] at (axis cs:11.5,3.5) {$a_0=3$};
      \end{axis}
    \end{tikzpicture}
  \end{center}

  \item Estimer graphiquement la valeur de $a_0$.
  \begin{tcolorbox}[corrigebox]
    Sur une période, le signal forme un triangle de hauteur $6$ sur la moitié de la période, puis vaut $0$ sur l'autre moitié. L'aire du triangle vaut $\frac{1}{2}\times 3 \times 6 = 9$. Rapportée à la période $T=6$, la valeur moyenne est $\frac{9}{6} = 1{,}5$...

    En fait, estimons plus soigneusement : le signal monte de $0$ à $6$ sur $[0;3]$, donc la valeur moyenne sur cet intervalle est $3$. Sur $[3;6]$ le signal vaut $0$. La moyenne globale est donc environ $\frac{3+0}{2} = 1{,}5$.

    \medskip
    \textit{Remarque :} on peut aussi estimer visuellement que la droite horizontale qui « équilibre les aires » se situe autour de $1{,}5$.
  \end{tcolorbox}

  \item Décomposer l'intégrale sur les deux intervalles et calculer $a_0$.
  \begin{tcolorbox}[corrigebox]
    \begin{align*}
      a_0 &= \frac{1}{T}\int_0^T s(t)\,dt = \frac{1}{6}\left[\int_0^3 2t\,dt + \int_3^6 0\,dt\right] \\
          &= \frac{1}{6}\left[\Big[t^2\Big]_0^3 + 0\right] \\
          &= \frac{1}{6}\left[9 - 0\right]
    \end{align*}
    \[\boxed{a_0 = \frac{9}{6} = \frac{3}{2} = 1{,}5}\]
  \end{tcolorbox}

  \item Quelle est l'interprétation physique de cette valeur si $s(t)$ représente une tension en volts ?
  \begin{tcolorbox}[corrigebox]
    Si $s(t)$ représente une tension, alors $a_0 = 1{,}5$\,V correspond à la \textbf{composante continue} (ou valeur moyenne) de cette tension. C'est la tension qu'indiquerait un voltmètre en mode DC branché aux bornes du signal. Autrement dit, un signal constant de $1{,}5$\,V transporterait la même énergie moyenne que ce signal périodique (en termes de valeur moyenne).
  \end{tcolorbox}
\end{enumerate}

\end{document}
